% Hand-written. Displays are tagged, not numbered, so that the equation
% numbers of the rest of the paper do not move.
\subsection{The variational modulus in closed form, and the critical window}
\label{sec:psi-window}
This subsection compares the variational modulus $\Psi$ of
\eqref{eq:v5-psi} with the closed quantile functional, and then evaluates
that functional for power laws with constants that do not depend on the
exponent. Its arguments are deterministic; the statistical conclusions
come from Corollary~\ref{fg:cor:general-laws} and
\cref{thm:entropy-confidence-explicit}.

Fix a probability law $\mu$ on $(0,1/2]$, $0<\rho\le1/4$ and $h>0$, and
put $B=h^{-2}$. Let $t$, $\nu$, $\kappa$ and $T_h$ be as in
\cref{fg:sec:comparison}, and write
$A_\sigma=\int d^{-2}\,\sigma(dd)$ for a measure $\sigma\le\mu$. Since
$\kappa$ has mass $1-\rho$ and lives on $[t,1/2]$, it occupies finitely
many bands and $0<A_\kappa<\infty$. Hence
$\Gamma_\kappa=A_\kappa(1+E_\kappa)$, where
\[
 E_\kappa=\Ent\bigl((a_j(\kappa)/A_\kappa)_j\bigr)
\]
is the entropy of the normalized band costs of the core. Define
\begin{align*}
 \Q^0(B)&=\min\bigl\{\rho,\ T_h+h\sqrt{\rho A_\kappa}\bigr\},\\
 \mathcal K^0(B)&=\min\Bigl\{\rho,\ \inf_{0\le\sigma\le\mu}
 \Bigl[h\sqrt{\rho A_\sigma}
       +\int\min\{1,h/d\}\,(\mu-\sigma)(dd)\Bigr]\Bigr\}.
 \tag{W1}\label{eq:psi-window-defs}
\end{align*}
These are $\Q(B)$ and $\mathcal K(B)$ with the cost entropy $\Gamma_\sigma$
replaced by the inverse-square cost $A_\sigma$. Because
$\Gamma_\sigma\ge A_\sigma$, we have $\mathcal K^0\le\mathcal K$ and
$\Q^0\le\Q$, and the choice $\sigma=\kappa$ gives $\mathcal K^0\le\Q^0$.

\begin{proposition}[$\Psi$ is the quantile form without its entropy]
\label{prop:psi-closed}
For every probability law $\mu$ on $(0,1/2]$, every $0<\rho\le1/4$ and
every $h>0$, with $B=h^{-2}$,
\[
 \frac{1}{\sqrt2}\,\Q^0(B)\le\Psi_{\mu,\rho}(h)\le\mathcal K^0(B)\le\Q^0(B).
 \tag{W2}\label{eq:psi-window-chain}
\]
Neither outer constant can be improved. If $h\ge t$ then
$\Psi_{\mu,\rho}(h)=\Q^0(B)=\rho$. If
$\mu=\rho\,\delta_{d_1}+(1-\rho)\,\delta_{d_2}$ with
$d_2=d_1\sqrt{(1-\rho)/\rho}\le1/2$ and $h\le d_1/2$, then
$\Psi_{\mu,\rho}(h)=\sqrt2\,h\rho/d_1=\Q^0(B)/\sqrt2$.
\end{proposition}

\begin{proof}
\emph{Upper bounds.} Let $0\le v\le1$ with $\int v\,d\mu\le\rho$, let
$0\le\sigma\le\mu$, and put $f=d\sigma/d\mu$. Since $v\le1$, the
integrand $\min\{v,h\sqrt v/d\}$ is at most $\min\{1,h/d\}$ and at most
$h\sqrt v/d$. Using the first bound against $1-f$, the second against
$f$, and the Cauchy--Schwarz inequality,
\begin{align*}
 \int\min\{v,h\sqrt v/d\}\,d\mu
 &\le\int\min\{1,h/d\}\,d(\mu-\sigma)
   +h\Bigl(\int fv\,d\mu\Bigr)^{1/2}\Bigl(\int fd^{-2}\,d\mu\Bigr)^{1/2}\\
 &\le\int\min\{1,h/d\}\,d(\mu-\sigma)+h\sqrt{\rho A_\sigma}.
\end{align*}
The left side is also at most $\int v\,d\mu\le\rho$. Taking the supremum
over $v$ and the infimum over $\sigma$ gives
$\Psi_{\mu,\rho}(h)\le\mathcal K^0(B)$.

\emph{Lower bound.} The function
$b(c)=\int\min\{1,c^2/d^2\}\,\mu(dd)$ is continuous and nondecreasing in
$c>0$ by dominated convergence, tends to zero as $c\downarrow0$, and
equals one at $c=1/2$. Choose $c$ with $b(c)=\rho$. Then
$v(d)=\min\{1,c^2/d^2\}$ is admissible in \eqref{eq:v5-psi}; let $V$ be
its objective value, so that $\Psi_{\mu,\rho}(h)\ge V$.

We first show $c<t$. If $c\ge t$, then $\mu((0,c])\ge\mu((0,t])\ge\rho$,
and $b(c)=\mu((0,c])+c^2\int_{(c,1/2]}d^{-2}\,d\mu=\rho$ forces
$\mu((c,1/2])=0$. Then $\mu((0,c])=1>\rho$, a contradiction. Hence
$\nu=\mu$ on $(0,c]$ and $\kappa((0,c])=0$.

If $h\ge c$, the objective integrand equals $v$ at every gap, so
$V=b(c)=\rho\ge\Q^0(B)$. Suppose $h<c$. Put
\begin{gather*}
 F=\mu((0,c]),\qquad C=\int_{(0,c]}\min\{1,h/d\}\,d\mu,\qquad
 m=\nu((c,t])=\rho-F,\\
 I_1=\int_{(c,t]}d^{-1}\,d\nu,\qquad I_2=\int_{(c,t]}d^{-2}\,d\nu .
\end{gather*}
On $(c,1/2]$ the integrand equals $hc/d^2$, and $\mu=\nu+\kappa$ there,
so $V=C+hc(I_2+A_\kappa)$. The budget identity $b(c)=\rho$ reads
$m=c^2(I_2+A_\kappa)$. Therefore
\[
 V=C+\frac{hm}{c},\qquad I_2+A_\kappa=\frac{m}{c^2},\qquad
 C\ge\frac{hF}{c}.
\]
The last inequality holds because $\min\{1,h/d\}\ge h/c$ for
$d\le c$. Moreover $T_h=C+hI_1$ and $I_1\le\sqrt{mI_2}$. By the
Cauchy--Schwarz inequality for the vectors
$(\sqrt m,\sqrt{F+m})$ and $(\sqrt{I_2},\sqrt{A_\kappa})$,
\[
 T_h+h\sqrt{\rho A_\kappa}
 \le C+h\bigl[\sqrt{mI_2}+\sqrt{(F+m)A_\kappa}\bigr]
 \le C+\frac{h}{c}\sqrt{m(F+2m)} .
\]
Write $u=F/c$ and $w=m/c$, so that $C\ge hu$ and $V=C+hw$. Then
\[
 \bigl[(\sqrt2-1)u+\sqrt2\,w\bigr]^2
 =(3-2\sqrt2)u^2+(4-2\sqrt2)uw+2w^2\ge uw+2w^2 ,
\]
that is, $\sqrt{w(u+2w)}\le(\sqrt2-1)u+\sqrt2\,w$. Consequently
\[
 T_h+h\sqrt{\rho A_\kappa}\le C+(\sqrt2-1)hu+\sqrt2\,hw
 \le\sqrt2\,(C+hw)=\sqrt2\,V .
\]
In both cases $\Q^0(B)\le\sqrt2\,V\le\sqrt2\,\Psi_{\mu,\rho}(h)$.

\emph{Sharpness.} If $h\ge t$, then $T_h=\rho$, so $\Q^0(B)=\rho$. Also
$h>c$, so the first case above gives $\Psi_{\mu,\rho}(h)\ge V=\rho$, and
$\Psi_{\mu,\rho}(h)\le\rho$ always. For the two-atom law, $d_1<d_2$, so
$\nu=\rho\,\delta_{d_1}$, $\kappa=(1-\rho)\,\delta_{d_2}$ and
$A_\kappa=\rho/d_1^2$; hence $\Q^0(B)=\min\{\rho,2h\rho/d_1\}=2h\rho/d_1$.
For any admissible $v$, the Cauchy--Schwarz inequality gives
\[
 \int\min\{v,h\sqrt v/d\}\,d\mu
 \le h\Bigl(\int v\,d\mu\Bigr)^{1/2}
      \Bigl(\frac{\rho}{d_1^2}+\frac{1-\rho}{d_2^2}\Bigr)^{1/2}
 \le\frac{\sqrt2\,h\rho}{d_1},
\]
and the lower bound just proved gives the reverse inequality.
\end{proof}

The proposition places $\Psi$ inside the framework of \cref{sec:entropy-tail}:
within a factor $\sqrt2$, it is the selection functional in which a
selected submeasure pays its inverse-square cost $A_\sigma$. The true rate
replaces $A_\sigma$ by $\Gamma_\sigma=A_\sigma(1+\Ent)$. Two consequences
follow at once.

\begin{corollary}[The $\Psi$ lower bound from the entropy--tail law]
\label{cor:psi-from-law}
For every probability law $\mu$ on $(0,1/2]$, every $0<\rho\le1/4$, every
$B>0$ and every $0<\delta\le1/4$,
\[
 \Psi_{\mu,\rho}(B^{-1/2})\le\mathcal K_{\mu,\rho}(B)\le\Q(B),
\]
and therefore
\[
 R_{\mu,\rho}(B)\ge10^{-17}\,\Psi_{\mu,\rho}(B^{-1/2}),\qquad
 r_{\rm exp}(B,\delta)\ge10^{-17}\,
 \Psi_{\mu,\rho}\Bigl(\sqrt{\log(1/\delta)/B}\Bigr).
 \tag{W3}\label{eq:psi-window-lower}
\]
\end{corollary}

\begin{proof}
The first display is \eqref{eq:psi-window-chain} with
$\mathcal K^0\le\mathcal K\le\Q$. The risk bound then follows from the
lower half of \eqref{eq:explicit-main-constants}. For the radius, apply
\cref{thm:entropy-confidence-explicit} at $\beta=B/\log(1/\delta)$ and use
$\Psi_{\mu,\rho}(\beta^{-1/2})\le\Q(\beta)$.
\end{proof}

The bounds \eqref{eq:psi-window-lower} are those of
Theorem~\ref{loc:thm:PsiLower}, now with an explicit constant and for
every $B>0$. This derivation uses neither the published single-coin
input of Lemma~\ref{loc:lem:LV} nor the hybrid argument of
\cref{app:localization}: it needs only the statements of the two cited
theorems and the inequality $\Psi\le\mathcal K$. For the risk alone,
\eqref{eq:psi-budget-lower} gives the better constant $1/8192$.

\begin{theorem}[How far $\Psi$ is from the rate]\label{thm:psi-gap}
For every probability law $\mu$ on $(0,1/2]$, every $0<\rho\le1/4$ and
every $B>0$, with $h=B^{-1/2}$,
\[
 \Psi_{\mu,\rho}(h)\le\mathcal K(B)\le\Q(B)
 \le\sqrt{1+E_\kappa}\;\Q^0(B)
 \le\sqrt{2(1+E_\kappa)}\;\Psi_{\mu,\rho}(h).
 \tag{W4}\label{eq:psi-window-gap}
\]
Consequently
\[
 10^{-17}\,\Psi_{\mu,\rho}(h)\le R_{\mu,\rho}(B)
 \le33000\sqrt{2(1+E_\kappa)}\;\Psi_{\mu,\rho}(h).
\]
If moreover $T_h\le h\sqrt{\rho A_\kappa}$ and
$T_h+h\sqrt{\rho\Gamma_\kappa}\le\rho$, then
$\Q(B)\ge\tfrac12\sqrt{1+E_\kappa}\,\Q^0(B)$, and so
\[
 R_{\mu,\rho}(B)\ge\frac{10^{-17}}{2(1+2\sqrt2)}
 \sqrt{1+E_\kappa}\;\Psi_{\mu,\rho}(h).
\]
The entropy satisfies $E_\kappa\le\log n_\kappa\le\log\log_2(1/t)$, where
$n_\kappa$ is the number of bands that $\kappa$ occupies.
\end{theorem}

\begin{proof}
The first two inequalities of \eqref{eq:psi-window-gap} are
Corollary~\ref{cor:psi-from-law}. For the third,
$T_h+h\sqrt{\rho A_\kappa(1+E_\kappa)}
 \le\sqrt{1+E_\kappa}\,(T_h+h\sqrt{\rho A_\kappa})$, and
$\min\{\rho,sx\}\le s\min\{\rho,x\}$ for $s\ge1$. The fourth is
\eqref{eq:psi-window-chain}. The risk bounds combine these with
\eqref{eq:explicit-main-constants}.

Under the two additional hypotheses,
$\Q(B)=T_h+h\sqrt{\rho A_\kappa(1+E_\kappa)}
 \ge\tfrac12\sqrt{1+E_\kappa}\,(T_h+h\sqrt{\rho A_\kappa})
 \ge\tfrac12\sqrt{1+E_\kappa}\,\Q^0(B)$,
and $\Q^0\ge\Psi$. Apply the lower bound of
Corollary~\ref{fg:cor:general-laws}.

An entropy is at most the logarithm of the size of its support. The bands
meeting $[t,1/2]$ are $I_0,\ldots,I_{j_0}$ with $t\in I_{j_0}$, and
$t\le b_{j_0}=2^{-(j_0+1)}$ gives $j_0+1\le\log_2(1/t)$.
\end{proof}

Thus $\Psi$ gives the rate, up to universal constants, on every family of
laws, rarities and budgets along which the core entropy $E_\kappa$ stays
bounded, for instance whenever the quantile core occupies a bounded
number of bands. When the core term dominates and the cap is inactive,
$\Psi$ falls short by a factor within universal constants of
$\sqrt{1+E_\kappa}$. For the
quadratic law of \cref{thm:v5-critical} the core occupies order $L$ bands
of comparable cost, so $1+E_\kappa\asymp1+\log L$; this is the factor by
which that theorem exceeds $\Psi$.

\subsubsection{Power laws, uniformly in the exponent}
Let $\mu_{\alpha,a}$ have distribution function $(d/a)^\alpha$ on $(0,a]$,
with $0<\alpha\le4$ and $0<a\le1/2$. For $0<\rho<1$ write
$s_\rho=\rho^{1/\alpha}$. The quantile is $t=as_\rho$, there are no atoms,
$\nu=\mu|_{(0,t]}$ and $\kappa=\mu|_{(t,a]}$. Define
\begin{align*}
 \Phi_\alpha(\rho)&=\alpha\int_{s_\rho}^1s^{\alpha-3}\,ds
 =\begin{cases}
   \dfrac{\alpha\,(1-\rho^{1-2/\alpha})}{\alpha-2},&\alpha\ne2,\\[1.5ex]
   \log(1/\rho),&\alpha=2,
  \end{cases}\\
 N_\alpha(\rho)&=\min\Bigl\{1+\frac{\log_2(1/\rho)}{\alpha},\
 \max\Bigl\{1,\frac1{|\alpha-2|}\Bigr\}\Bigr\},
 \tag{W5}\label{eq:power-phi-n}
\end{align*}
with $1/0=\infty$, so that $N_2(\rho)=1+\tfrac12\log_2(1/\rho)$. The
function $\Phi_\alpha$ is continuous in $\alpha$. The number $N_\alpha(\rho)$
counts the dyadic scales of the core that carry comparable cost: all of
them when $|\alpha-2|\log(1/\rho)$ is small, and about $1/|\alpha-2|$ of
them otherwise.

\begin{lemma}[The quantile core of a power law]\label{lem:power-core}
For $0<\alpha\le4$, $0<a\le1/2$ and $0<\rho<1$,
\[
 a^2A_\kappa=\Phi_\alpha(\rho),\qquad
 \tfrac12\bigl[1+\log N_\alpha(\rho)\bigr]\le1+E_\kappa
 \le3.83\bigl[1+\log N_\alpha(\rho)\bigr].
 \tag{W6}
\]
\end{lemma}

\begin{proof}
The substitution $d=as$ gives
$A_\kappa=\alpha a^{-\alpha}\int_t^ad^{\alpha-3}\,dd=a^{-2}\Phi_\alpha(\rho)$.

Parametrize the core by $y=\log_2(a/d)\in[0,M)$, where
$M=\log_2(1/s_\rho)=\log_2(1/\rho)/\alpha$. The cost measure
$d^{-2}\kappa(dd)$ becomes
$\alpha(\log2)a^{-2}\,2^{-ey}\,dy$ with $e=\alpha-2\in(-2,2]$, and the
band $I_j$ becomes the unit interval
$[j+1+\log_2a,\ j+2+\log_2a)$. Hence $E_\kappa$ is the entropy of the
partition of $[0,M)$ into these unit cells under the probability measure
$P$ whose density is proportional to $2^{-ey}$. Write $p_k$ for the cell
probabilities.

\emph{Upper bound.} At most $\lceil M\rceil+1\le M+2$ cells meet
$[0,M)$, so $E_\kappa\le\log(M+2)\le\log2+\log(1+M)$. Now let $e\ne0$.
Replacing $y$ by $M-y$ if $e<0$, we may assume that $P$ has density
proportional to $2^{-|e|y}$ on $[0,M]$; the cells remain unit intervals.
Number them $k=0,1,2,\ldots$ in increasing order of $y$. Cell $k\ge1$
lies in $[k-1,\infty)$, and
$P([s,\infty))=(2^{-|e|s}-2^{-|e|M})/(1-2^{-|e|M})\le\theta^s$ for
$0\le s\le M$, where $\theta=2^{-|e|}\in[1/4,1)$. Therefore
$\sum_kkp_k=\sum_{k\ge1}P(\text{cells }k,k+1,\ldots)
 \le\sum_{k\ge1}\theta^{k-1}=1/(1-\theta)$.
Gibbs' inequality with the reference probabilities $(1-\theta)\theta^k$
gives
\[
 E_\kappa\le\log\frac1{1-\theta}+\log\frac1\theta\sum_kkp_k
 \le\log\frac1{1-\theta}+\frac{\log(1/\theta)}{1-\theta}.
\]
The last fraction decreases in $\theta$, so it is at most
$\tfrac43\log4<1.849$. Concavity of $z\mapsto1-2^{-z}$ on $[0,2]$ gives
$1-\theta\ge\tfrac38|e|$, so
$E_\kappa<\log\tfrac83+1.849+\log(1/|e|)<2.83+\log(1/|e|)$, and
$E_\kappa<2.83$ if $|e|\ge1$.
Together,
\[
 1+E_\kappa\le\min\bigl\{1+\log2+\log(1+M),\
 3.83+\log\max\{1,1/|e|\}\bigr\}\le3.83\,[1+\log N_\alpha(\rho)] .
\]

\emph{Lower bound.} An entropy is at least $\log(1/\max_kp_k)$. A cell
has length at most one, so $\max_kp_k$ is at most the supremum of the
density of $P$, which is $M^{-1}\zeta/(1-e^{-\zeta})$ with
$\zeta=|e|M\log2$, and $M^{-1}$ if $e=0$. Since
$\zeta/(1-e^{-\zeta})\le1+\zeta$,
\[
 \max_kp_k\le\min\Bigl\{1,\frac1M+|e|\log2\Bigr\}
 \le\min\Bigl\{1,\frac{1+\log2}{z}\Bigr\},\qquad
 z=\min\{M,1/|e|\}.
\]
Hence $1+E_\kappa\ge1+\log\max\{1,z/(1+\log2)\}$. The right side is at
least $\tfrac12[1+\log(1+z)]$: for $z\le1+\log2$ this reads
$1\ge\tfrac12[1+\log(2+\log2)]$, and for $z\ge1+\log2$ it follows from
$\log(1+z)\le\log z+\log\bigl(1+1/(1+\log2)\bigr)$ and $\log z\ge0.52$.
Finally $N_\alpha(\rho)\le1+z$.
\end{proof}

For $0<x<s_\rho$ put
\[
 T_\alpha(x,\rho)=x^\alpha+\alpha x\int_x^{s_\rho}s^{\alpha-2}\,ds,
 \qquad
 Y_\alpha(x,\rho)=x\sqrt{\rho\,\Phi_\alpha(\rho)\,[1+\log N_\alpha(\rho)]},
\]
and define
\[
 \mathcal R_\alpha(x,\rho)=
 \begin{cases}
  \min\{\rho,\ T_\alpha(x,\rho)+Y_\alpha(x,\rho)\},&0<x<s_\rho,\\
  \rho,&x\ge s_\rho .
 \end{cases}
 \tag{W7}\label{eq:power-uniform-rate}
\]

\begin{theorem}[Power margins, uniformly in the exponent]
\label{thm:power-uniform}
Let $0<\alpha\le4$, $0<a\le1/2$, $0<\rho\le1/4$ and $B>0$, and put
$x=(a\sqrt B)^{-1}$. For the law $\mu_{\alpha,a}$,
\[
 \frac1{\sqrt2}\,\mathcal R_\alpha(x,\rho)\le\Q(B)
 \le2\,\mathcal R_\alpha(x,\rho),
 \qquad
 10^{-18}\,\mathcal R_\alpha(x,\rho)\le R_{\mu_{\alpha,a},\rho}(B)
 \le66000\,\mathcal R_\alpha(x,\rho).
 \tag{W8}\label{eq:power-uniform}
\]
For $0<\delta\le1/4$ and $x_\delta=a^{-1}\sqrt{\log(1/\delta)/B}$, the
confidence radii of \cref{app:known-confidence} satisfy
\[
 10^{-18}\,\mathcal R_\alpha(x_\delta,\rho)\le r_{\rm exp}(B,\delta)
 \le r_{\rm hard}(B,\delta)
 \le1.5\cdot10^{6}\,\mathcal R_\alpha(x_\delta,\rho).
\]
The constants do not depend on $\alpha$, $a$, $\rho$, $B$ or $\delta$.
\end{theorem}

\begin{proof}
Let $h=B^{-1/2}=ax$. If $x\ge s_\rho$ then $h\ge t$ and
$\Q(B)=\rho=\mathcal R_\alpha(x,\rho)$. If $x<s_\rho$, the substitution
$d=as$ gives
$T_h=\mu((0,h])+h\int_{(h,t]}d^{-1}\,\mu(dd)=T_\alpha(x,\rho)$, and
Lemma~\ref{lem:power-core} gives
$h\sqrt{\rho\Gamma_\kappa}=x\sqrt{\rho\,\Phi_\alpha(\rho)(1+E_\kappa)}$
with $1/\sqrt2\le\sqrt{(1+E_\kappa)/(1+\log N_\alpha(\rho))}\le\sqrt{3.83}<2$.
This proves the bounds on $\Q(B)$. The risk bounds follow from
Corollary~\ref{fg:cor:general-laws}, since
$10^{-17}/[(1+2\sqrt2)\sqrt2]>10^{-18}$. For the radii, apply
\cref{thm:entropy-confidence-explicit} with
$\mathcal K(\beta)\le\Q(\beta)$ at $\beta=B/\log(1/\delta)$, for which
$(a\sqrt\beta)^{-1}=x_\delta$.
\end{proof}

\begin{corollary}[The critical window]\label{cor:critical-window}
Let $3/2\le\alpha\le3$, with the other parameters as in
\cref{thm:power-uniform}, and write $Y=Y_\alpha(x,\rho)$. Then
$\min\{\rho,Y\}\le\mathcal R_\alpha(x,\rho)\le5.13\min\{\rho,Y\}$, and
therefore
\[
 10^{-18}\min\{\rho,Y\}\le R_{\mu_{\alpha,a},\rho}(B)
 \le3.4\cdot10^{5}\min\{\rho,Y\},
 \qquad
 Y=\frac{\sqrt{\rho\,\Phi_\alpha(\rho)\,[1+\log N_\alpha(\rho)]}}{a\sqrt B}.
 \tag{W9}\label{eq:psi-window-critical}
\]
Moreover
$\tfrac1{\sqrt2}\min\{\rho,Y^0\}\le\Psi_{\mu,\rho}(B^{-1/2})
 \le5.13\min\{\rho,Y^0\}$ with
$Y^0=x\sqrt{\rho\,\Phi_\alpha(\rho)}$. In particular, whenever $Y\le\rho$,
the ratio $R_{\mu_{\alpha,a},\rho}(B)/\Psi_{\mu,\rho}(B^{-1/2})$ is
within universal constants of $\sqrt{1+\log N_\alpha(\rho)}$.
\end{corollary}

\begin{proof}
For $\alpha>1$ and $x<s_\rho$,
$T_\alpha(x,\rho)=x[\alpha s_\rho^{\alpha-1}-x^{\alpha-1}]/(\alpha-1)
 \le\frac{\alpha}{\alpha-1}\,x\rho^{1-1/\alpha}$, and
$\alpha/(\alpha-1)\le3$ for $\alpha\ge3/2$, $\le2$ for $\alpha\ge2$.

If $3/2\le\alpha\le2$, then $s_\rho\le1/2$, and since $s^{\alpha-3}$
decreases,
$\Phi_\alpha(\rho)\ge\alpha\int_{s_\rho}^{2s_\rho}s^{\alpha-3}\,ds
 \ge\alpha2^{\alpha-3}\rho^{1-2/\alpha}\ge0.53\,\rho^{1-2/\alpha}$.
Hence $\sqrt{\rho\Phi_\alpha(\rho)}\ge0.728\,\rho^{1-1/\alpha}$. If
$2\le\alpha\le3$, then $s^{\alpha-3}\ge1$ on $(0,1]$, so
$\Phi_\alpha(\rho)\ge\alpha(1-s_\rho)\ge2(1-4^{-1/3})>0.74$ and
$\sqrt{\rho\Phi_\alpha(\rho)}\ge0.86\sqrt\rho\ge0.86\,\rho^{1-1/\alpha}$.
In both cases $T_\alpha(x,\rho)\le4.13\,Y^0\le4.13\,Y$ for $x<s_\rho$,
which gives $\min\{\rho,Y\}\le\mathcal R_\alpha\le5.13\min\{\rho,Y\}$ there.
For $x\ge s_\rho$ the same lower bounds give
$Y\ge Y^0\ge s_\rho\sqrt{\rho\Phi_\alpha(\rho)}\ge0.728\,\rho$, using
$\rho^{1/\alpha+1/2}\ge\rho$ when $\alpha\ge2$; hence
$\mathcal R_\alpha=\rho\le1.38\min\{\rho,Y\}$. The risk bounds follow
from \cref{thm:power-uniform}.

The same argument with $E_\kappa$ deleted gives
$\min\{\rho,Y^0\}\le\Q^0(B)\le5.13\min\{\rho,Y^0\}$, and
Proposition~\ref{prop:psi-closed} transfers this to $\Psi$. If $Y\le\rho$
then also $Y^0\le\rho$, and $Y/Y^0=\sqrt{1+\log N_\alpha(\rho)}$.
\end{proof}

\begin{remark}[Reading the uniform rate]\label{rem:window}
At $\alpha=2$ and $a=1/4$, \eqref{eq:psi-window-critical} is the rate of
\cref{thm:v5-critical}: $\Phi_2(\rho)=\log(1/\rho)$ and
$N_2(\rho)=1+\tfrac12\log_2(1/\rho)$ are within universal factors of $L$.

Away from two the formula gives the rows of \cref{sec:rare} together
with the dependence of their constants on the exponent. Put
$u=|1-2/\alpha|\log(1/\rho)$ and suppose $u\ge1$, with
$3/2\le\alpha\le3$ and $\alpha\ne2$. Then
$N_\alpha(\rho)=1/|\alpha-2|$, and $\Phi_\alpha(\rho)$ lies between
$1-e^{-1}$ and $1$ times $\alpha/(\alpha-2)$ for $\alpha>2$, and times
$\alpha\rho^{1-2/\alpha}/(2-\alpha)$ for $\alpha<2$. Hence, with
$\epsilon=|\alpha-2|$,
\[
 R_{\mu_{\alpha,a},\rho}(B)\asymp
 \min\Bigl\{\rho,\ x\,\rho^{\min\{1/2,\,1-1/\alpha\}}
 \sqrt{\frac{1+\log(1/\epsilon)}{\epsilon}}\Bigr\}
\]
with universal constants. The constants of the fixed-exponent rows
therefore grow like $\sqrt{\log(1/\epsilon)/\epsilon}$ as the exponent
approaches two, and they stop growing once $\epsilon\log(1/\rho)\le1$:
at rarity $\rho$ the exponents within $1/\log(1/\rho)$ of two behave like
the quadratic law, with $\Phi_\alpha(\rho)$ and $N_\alpha(\rho)$ both of
order $\log(1/\rho)$.
\end{remark}

\paragraph{Status.}
The statements of this subsection have paper proofs. They use
Corollary~\ref{fg:cor:general-laws} and
\cref{thm:entropy-confidence-explicit}, whose scopes are recorded in
\cref{sec:verification}, and elementary inequalities that have not been
formalized. The floating-point script
\path{research/2026-11-09-psi-closed-form/check.py} tests every inequality
and constant above on random finite laws and on a grid of power laws; it
is numerical evidence only.
